\documentclass[11pt]{article}

\usepackage[margin=1in]{geometry}
\usepackage{graphicx}
\usepackage{amsmath,amssymb}
\usepackage{booktabs}
\usepackage{microtype}
\usepackage{hyperref}
\usepackage{cleveref}
\usepackage{natbib}

\title{A TOGA-inspired Method for ncRNA Analysis}
\author{Author One$^{1}$, Author Two$^{2}$\\
$^{1}$Affiliation One\\
$^{2}$Affiliation Two}
\date{\today}

\begin{document}
\maketitle

\begin{abstract}
Placeholder abstract. Summarize the problem, the method (TOGA-like), key results, and why it matters for ncRNAs.
\end{abstract}

\section{Introduction}
Motivate ncRNA analysis, gaps in current methods, and why a TOGA-style approach helps. Outline contributions.

\section{Related Work}
Summarize relevant ncRNA methods and any TOGA lineage or inspiration. Add citations.

\section{Method}
Describe the pipeline end-to-end.
\subsection{Overview}
High-level diagram and inputs/outputs.
\subsection{Data and Preprocessing}
Datasets, filtering, quality control.
\subsection{Model/Algorithm}
Core method details, notation, and rationale.
\subsection{Implementation}
Key engineering details and runtime considerations.

\section{Experiments}
\subsection{Setup}
Train/test splits, baselines, metrics.
\subsection{Results}
Tables/figures with quantitative results.
\subsection{Ablations}
Component-wise analysis.

\section{Discussion}
Interpret results, limitations, and biological insights.

\section{Future Work}
Planned extensions and open questions.

\section*{Acknowledgments}
Optional.

\bibliographystyle{plainnat}
\bibliography{references}

\end{document}
